\section*{Chapitre \ref{ch:b2:diatomic-mos} --- Orbitales moléculaires des molécules diatomiques}

\begin{solution}{exo:b2:diatomic-mos:1}
$E_+ = (-13.6 - 9.0)/1.60 = \qty{-14.1}{eV}$~; $E_- = (-13.6 + 9.0)/0.40 = \qty{-11.5}{eV}$.
Deux électrons dans $\psi_+$~: $-28.25$ contre $2 \times (-13.6) = \qty{-27.2}{eV}$, soit une liaison de
\qty{1.05}{eV} dans ce modèle grossier.
\end{solution}

\begin{solution}{exo:b2:diatomic-mos:2}
\ce{He2}~: $(2 - 2)/2 = 0$~; \ce{He2+}~: $(2 - 1)/2 = 0.5$~; \ce{Li2}~: 1~; \ce{B2}~: six électrons de
valence, $\sigma_{2s}^2\sigma_{2s}^{*2}\pi_{2p}^2$, \omterm{def:b2:diatomic-mos:bond-order}{indice de liaison} 1.
\end{solution}

\begin{solution}{exo:b2:diatomic-mos:3}
\ce{B2} (deux électrons célibataires dans les deux orbitales $\pi_{2p}$, avec interaction s--p) et
\ce{O2} (deux dans $\pi^*$). \ce{C2}, \ce{N2}, \ce{F2} sont \omterm{def:b2:diatomic-mos:magnetism}{diamagnétiques}.
\end{solution}

\begin{solution}{exo:b2:diatomic-mos:4}
Non nul~: (1s, $2\mathrm p_z$), ($2\mathrm p_x$, $2\mathrm p_x$), (2s, $2\mathrm p_z$).
Nul par symétrie~: (1s, $2\mathrm p_x$) et ($2\mathrm p_x$, $2\mathrm p_y$).
\end{solution}

\begin{solution}{exo:b2:diatomic-mos:5}
L'ionisation retire un électron liant~: \omterm{def:b2:diatomic-mos:bond-order}{indice de liaison} 2.5 au lieu de 3, donc une liaison plus longue et
plus faible, comme le montre $D_0(\ce{N2+}) = \qty{8.71}{eV} < D_0(\ce{N2}) = \qty{9.76}{eV}$.
\end{solution}

\begin{solution}{exo:b2:diatomic-mos:6}
Les orbitales de l'oxygène sont plus basses, car l'oxygène est plus électronégatif. D'après
le \cref{thm:b2:diatomic-mos:heteronuclear}, les \omterm{def:b2:diatomic-mos:bonding}{orbitales liantes} sont développées surtout sur
l'atome le plus bas, l'oxygène, et les orbitales plus hautes sur le carbone. La plus haute orbitale
occupée, une orbitale $\sigma$ à caractère surtout carboné, orientée à l'opposé de l'oxygène, est
le doublet non liant du carbone.
\end{solution}

\begin{solution}{exo:b2:diatomic-mos:7}
Seule la $2\mathrm p_z$ du fluor (selon l'axe) a la symétrie de la 1s de l'hydrogène~:
elles forment une orbitale $\sigma$ développée surtout sur F et une $\sigma^*$ développée surtout sur H. Les
$2\mathrm p_x$, $2\mathrm p_y$ du fluor (recouvrement nul) et sa 2s basse (trop éloignée en énergie)
restent non liantes~: trois doublets non liants. Les huit électrons de valence remplissent 2s, $\sigma$
et les deux p non liantes.
\end{solution}

\begin{solution}{exo:b2:diatomic-mos:8}
$\bar\alpha = \qty{-12}{eV}$, $\Delta = \qty{4}{eV}$, $\sqrt{4 + 9} = \qty{3.61}{eV}$~:
$E = \qty{-15.6}{eV}$ et \qty{-8.4}{eV}. Pour le niveau inférieur, $(\alpha_B - E)c_B = -\beta
c_A$~: $1.61c_B = 3.0c_A$, $c_A/c_B = 0.535$, $c_B^2 = 1/(1 + 0.286) = 0.78$.
\end{solution}

\begin{solution}{exo:b2:diatomic-mos:9}
Huit électrons de valence~: $\sigma_{2s}^2\sigma_{2s}^{*2}\pi_{2p}^4$, l'orbitale $\sigma_{2p}$
(repoussée au-dessus de $\pi_{2p}$ par l'interaction) étant vide. \omterm{def:b2:diatomic-mos:bond-order}{Indice de liaison} $(6 - 2)/2 = 2$, dû aux
deux orbitales $\pi$ pleines.
\end{solution}

\begin{solution}{exo:b2:diatomic-mos:10}
NO, onze électrons de valence, dont un dans $\pi^*$~: \omterm{def:b2:diatomic-mos:bond-order}{indice de liaison} 2.5. \ce{NO+}, dix~: \omterm{def:b2:diatomic-mos:bond-order}{indice de liaison}
3. $D_0(\ce{NO+}) = 6.51 + 13.62 - 9.26 = \qty{10.87}{eV}$~: liaison bien plus forte que dans NO, puisque
l'électron retiré était antiliant.
\end{solution}

\begin{solution}{exo:b2:diatomic-mos:11}
$E_+ + E_- = \dfrac{(\alpha + \beta)(1 - S) + (\alpha - \beta)(1 + S)}{1 - S^2} = \dfrac{2(\alpha -
\beta S)}{1 - S^2}$. Cette somme dépasse $2\alpha$ lorsque $\alpha - \beta S > \alpha - \alpha S^2$, c'est-à-dire
$\beta < \alpha S$ (en divisant par $-S < 0$). Avec quatre électrons, deux dans chaque orbitale, l'énergie
totale est supérieure à celle des orbitales séparées~: deux orbitales pleines se repoussent.
\end{solution}

\begin{solution}{exo:b2:diatomic-mos:12}
\omterm{def:b2:diatomic-mos:bond-order}{Indices de liaison} 2.5, 2, 1.5, 1~: dans cet ordre, les liaisons s'allongent et s'affaiblissent. Le peroxyde
d'hydrogène contient la liaison simple O--O peroxyde (motif \ce{O2^2-})~; \ce{KO2} contient
l'ion superoxyde \ce{O2-}, \omterm{def:b2:diatomic-mos:magnetism}{paramagnétique}.
\end{solution}

\begin{solution}{pb:b2:diatomic-mos:1}
\textbf{1.} Douze.
\textbf{2.} $\sigma_{2s}^2\,\sigma_{2s}^{*2}\,\sigma_{2p}^2\,\pi_{2p}^4\,\pi_{2p}^{*2}$, les deux électrons $\pi^*$
étant dans des orbitales différentes.
\textbf{3.} $(8 - 4)/2 = 2$.
\textbf{4.} Deux~: \ce{O2} est \omterm{def:b2:diatomic-mos:magnetism}{paramagnétique}.
\textbf{5.} Il apparie tous les électrons, alors que deux d'entre eux sont seuls dans deux orbitales
$\pi^*$ dégénérées.
\textbf{6.} $D_0 = 2 \times 246.79 = \qty{493.6}{kJ/mol}$, \qty{5.12}{eV}.
\textbf{7.} \ce{O2+}~: $\pi^{*1}$, \omterm{def:b2:diatomic-mos:bond-order}{indice de liaison} 2.5.
\textbf{8.} \ce{O2-}~: $\pi^{*3}$, \omterm{def:b2:diatomic-mos:bond-order}{indice de liaison} 1.5.
\textbf{9.} \ce{O2^2-}~: $\pi^{*4}$, \omterm{def:b2:diatomic-mos:bond-order}{indice de liaison} 1.
\textbf{10.} \ce{O2+} un, \ce{O2-} un, \ce{O2^2-} aucun.
\textbf{11.} \ce{O2+} $<$ \ce{O2} $<$ \ce{O2-} $<$ \ce{O2^2-}.
\textbf{12.} L'ion peroxyde.
\textbf{13.} \ce{O2+}, \ce{O2} et \ce{O2-}.
\textbf{14.} D'une orbitale $\pi^*$.
\textbf{15.} L'orbitale $\pi^*$, antiliante, est au-dessus du niveau 2p de l'atome~: son
électron est moins lié.
\textbf{16.} Deux chemins mènent de \ce{O2} à $\ce{O} + \ce{O+} + \mathrm e^-$~: dissocier puis
ioniser un atome, $D_0(\ce{O2}) + I(\ce{O})$~; ou ioniser la molécule puis dissocier l'ion,
$I(\ce{O2}) + D_0(\ce{O2+})$. Ils sont égaux~: $D_0(\ce{O2+}) = 5.12 + 13.62 - 12.07 = \qty{6.66}{eV}$.
\textbf{17.} Supérieure à $D_0(\ce{O2}) = \qty{5.12}{eV}$~: retirer un électron antiliant
renforce la liaison (indice 2.5).
\textbf{18.} $D_0(\ce{N2}) = 2 \times 470.82/96.485 = \qty{9.76}{eV}$~; $D_0(\ce{N2+}) = 9.76 +
14.53 - 15.58 = \qty{8.71}{eV}$.
\textbf{19.} La plus haute orbitale occupée de \ce{N2} est liante, au-dessous du niveau 2p de
l'atome.
\textbf{20.} Lorsque l'électron retiré provient d'une \omterm{def:b2:diatomic-mos:bonding}{orbitale antiliante}.
\textbf{21.} La règle de Hund (nombre maximal de spins parallèles dans des orbitales dégénérées).
\textbf{22.} Les deux électrons $\pi^*$ dans la même orbitale, spins appariés.
\textbf{23.} Deux électrons dans la même orbitale se repoussent davantage, et perdent la
stabilisation d'échange des spins parallèles.
\textbf{24.} \omterm{def:b2:diatomic-mos:bond-order}{Indice de liaison} toujours 2~; tous les électrons appariés~: espèce \omterm{def:b2:diatomic-mos:magnetism}{diamagnétique}.
\textbf{25.} \ce{O2+} a un \omterm{def:b2:diatomic-mos:bond-order}{indice de liaison} $\boldsymbol{2.5}$ et une énergie de dissociation de
$\boldsymbol{\approx \qty{6.66}{eV}}$, contre 2 et \qty{5.12}{eV} pour \ce{O2}.
\end{solution}
